% T33 - Informatique embarquée, programmation et réseaux industriels - créée le 2026-11-10.
% Présence : 11/13 années 2013-2025 où un sujet mobilise algorithme, pseudo-code, Python, C,
% microcontrôleur, registre, UART/RS-485, I2C/SPI, Modbus, Ethernet ou IHM (hors le mot
% « programme », trop ambigu) - recherche sur les .md du corpus le 2026-11-10.
% QT1 : E1 2026 D10 (corrigé officiel RJ26 = question D11, Rapport_Jury_2026_3eConcours.pdf p.18).
% QT2 : E2 disciplinaire IE 2023 E4-E8 (corrigé officiel Rapport_Jury_2023.pdf p.12).
%       Le rapport écrit « 1,03 m » pour E8 : le calcul donne 1,03 cm (signalé dans la QT).
% QT3 : E2 disciplinaire IE 2025 Q4-Q5 (corrigé officiel Rapport_Jury_2025.pdf p.8).
% Citation Bloc 4 : Rapport_Jury_2018.md l.158 (avant-propos commun à toutes les options),
% verbatim vérifié ; coupure finale marquée […].
\expandafter\ifx\csname CAPETINCLUDEMODE\endcsname\relax
\documentclass[11pt,a4paper]{article}
% preamble.tex — packages et commandes communes a toutes les fiches CAPET SII IE
\input{glyphtounicode}
\pdfgentounicode=1
\usepackage{cmap}
\usepackage[T1]{fontenc}
\usepackage[utf8]{inputenc}
\usepackage[french]{babel}
% babel[french] active a \begin{document} des raccourcis de ponctuation (:;!?)
% qui cassent la syntaxe to[...] de circuitikz (erreur "+ or - expected") ;
% on les desactive, ce qui ne change que l'espacement fin avant ces signes.
\AtBeginDocument{\shorthandoff{:;!?}}
\usepackage[a4paper, top=1.6cm, bottom=1.6cm, left=1.4cm, right=1.4cm]{geometry}
\setlength{\headheight}{14.5pt}
\usepackage{amsmath, amssymb, amsfonts}
\usepackage{siunitx}
\sisetup{locale=FR, output-decimal-marker={,}, inter-unit-product=\ensuremath{{}\cdot{}},
         group-minimum-digits=4}
% Unites non fournies par siunitx, declarees ici une fois pour toutes (fiches + exercices).
% Les gardes \ifdefined evitent tout conflit avec une redeclaration locale restee dans un
% fichier : ne pas retirer les gardes, ne pas redeclarer ces unites dans les fichiers.
\DeclareSIUnit{\jour}{j}
\ifdefined\var\else\DeclareSIUnit{\var}{var}\fi   % puissance reactive
\ifdefined\vad\else\DeclareSIUnit{\vad}{VAD}\fi   % puissance deformante
\ifdefined\tr\else\DeclareSIUnit{\tr}{tr}\fi      % tour (tr/min)
\usepackage{array, booktabs, multirow, tabularx}
\usepackage{xcolor}
\usepackage[most, breakable]{tcolorbox}
\usepackage{multicol}
\usepackage{enumitem}
\usepackage{mathtools}
\usepackage{lmodern}
% Refonte visuelle 2026 (direction "moderne coloré") : Lato pour le texte,
% maths sans-serif assorties via newtxsf (charge apres amsmath/amssymb).
\usepackage[default]{lato}
\usepackage{newtxsf}
\usepackage{microtype}
\usepackage{fancyhdr}
\usepackage{tikz}
\usetikzlibrary{positioning,arrows.meta}
\usepackage{circuitikz}
\usepackage{lastpage}

% ── Couleurs par domaine (palette modernisée 2026) ────────────────────
\definecolor{colElec}{RGB}{29, 78, 216}        % Électrotechnique — bleu vif
\definecolor{colElecLight}{RGB}{219, 234, 254} % fond bleu clair
\definecolor{colAuto}{RGB}{21, 128, 61}        % Automatique — vert
\definecolor{colMeca}{RGB}{194, 65, 12}        % Mécanique — orange
\definecolor{colTherm}{RGB}{185, 28, 28}       % Thermique — rouge
\definecolor{colQT}{RGB}{232, 244, 233}        % (hérité, plus utilisé par questiontype)
\definecolor{colJury}{RGB}{255, 248, 225}      % fond jury (crème)
\definecolor{colJuryFrame}{RGB}{202, 138, 4}   % bordure/bandeau jury (ambre)
\definecolor{colGrayLight}{RGB}{245, 245, 247} % fond gris neutre
\definecolor{colSeparator}{RGB}{180, 180, 180} % gris séparateur
\definecolor{colExo}{RGB}{8, 145, 178}          % Banque d'exercices — cyan (distinct des fiches)
% ── Encre et gris de la refonte ───────────────────────────────────────
\definecolor{colInk}{RGB}{23, 32, 42}          % texte courant (quasi-noir)
\definecolor{colSlate}{RGB}{83, 98, 118}       % gris ardoise (méta, en-têtes)
\definecolor{colRule}{RGB}{203, 213, 225}      % filets clairs
\definecolor{colJuryBg}{RGB}{255, 248, 225}    % fond crème jury
\definecolor{colPresence}{RGB}{83, 98, 118}    % bloc 5 (présence dans les sujets)
\AtBeginDocument{\color{colInk}}

% ── Style encadré exercice + corrigé (banque d'exercices) ─────────────
\tcbset{
  exoBox/.style={
    breakable,
    enhanced,
    colback=colExo!8!white,
    colframe=colExo,
    colbacktitle=colExo,
    coltitle=white,
    fonttitle=\bfseries\small,
    left=4pt, right=4pt, top=3pt, bottom=3pt,
    before skip=5pt, after skip=5pt,
  }
}

% ── Bascule enonce / corrige (banque d'exercices) ─────────────────────
% \CAPETCORRIGE est defini en ligne de commande (pdflatex "\def\CAPETCORRIGE{1}\input{...}")
% pour compiler la variante corriges d'un meme fichier source ; par defaut (non defini),
% la compilation produit la variante enonces. \CAPETSHOW{enonce}{corrige} choisit la branche.
%
% TROIS MODES (refonte 2026-09-20, manuel de revision) :
%   - aucun drapeau   -> imprime l'enonce            (variante enonces, INCHANGE)
%   - \CAPETCORRIGE   -> imprime le corrige          (variante corriges, INCHANGE)
%   - \CAPETMANUEL    -> imprime l'enonce ET met le corrige de cote pour l'annexe
%
% Le mode manuel stocke chaque corrige dans une macro (\csgdef), sans ecriture
% fichier ni re-tokenisation : le contenu (circuitikz, accents, %) traverse intact.
% L'annexe les rejoue ensuite dans l'ordre via \CAPETANNEXE.
\usepackage{etoolbox}

\newcounter{CAPETexo}          % numero continu d'exercice dans le manuel
\setcounter{CAPETexo}{0}
\newcounter{CAPETannexe}       % compteur de relecture en annexe

% Titre courant de la notion en cours (pose par \CAPETNOTION), sert d'intertitre annexe
\newcommand{\CAPETnotioncode}{}
\newcommand{\CAPETnotiontitre}{}

% \CAPETNOTION{code}{titre} : ouvre une notion dans le manuel (memorise le libelle,
% et marque le rang du premier exercice de la notion pour l'intertitre d'annexe).
\newcommand{\CAPETNOTION}[2]{%
  \renewcommand{\CAPETnotioncode}{#1}%
  \renewcommand{\CAPETnotiontitre}{#2}%
  \csgdef{CAPETnotionstart\the\numexpr\value{CAPETexo}+1\relax}{#1 -- #2}%
}

% ── \degree robuste ───────────────────────────────────────────────────
% Certains fichiers d'exercices utilisent \degree. Ni le preambule ni siunitx
% ne le definissent, et \providecommand local ne suffit pas : sous hyperref,
% \degree existe deja comme commande factice (\ERROR), donc \providecommand
% s'abstient et l'appel echoue a la compilation du manuel (erreur invisible
% en compilation isolee). On le (re)definit donc inconditionnellement ici,
% apres le chargement des packages, pour tout le corpus.
\AtBeginDocument{\providecommand{\degree}{}\renewcommand{\degree}{\ensuremath{^\circ}}}

% Libelle decoratif : les fichiers d'exercices utilisent aussi \CAPETSHOW dans
% leur \fancyhead et leur bloc de titre, pour afficher "(enonces)" ou "(corriges)".
% Ces appels-la ne sont PAS des exercices : en mode manuel ils ne doivent ni
% consommer un numero, ni alimenter l'annexe. On les reconnait a leur premier
% argument, strictement egal au libelle ci-dessous.
\newcommand{\CAPETLIBELLEENONCES}{(\'enonc\'es)}

\ifdefined\CAPETMANUEL
  % --- Mode manuel : enonce en place, corrige differe vers l'annexe ---
  \newcommand{\CAPETSHOW}[2]{%
    \def\CAPETargun{#1}%
    \ifx\CAPETargun\CAPETLIBELLEENONCES
      #1%            appel decoratif : rendu tel quel, hors numerotation
    \else
      \CAPETSHOWEXO{#1}{#2}%
    \fi
  }
  \newcommand{\CAPETSHOWEXO}[2]{%
    \stepcounter{CAPETexo}%
    \csgdef{CAPETcorrige\theCAPETexo}{#2}%
    \edef\CAPETtmpnum{\theCAPETexo}%
    #1%
    \par\vspace{2pt}%
    {\footnotesize\itshape\color{colExo}%
     \hyperref[capet:cor:\CAPETtmpnum]{Corrige de l'exercice \CAPETtmpnum\ $\rightarrow$ p.~\pageref{capet:cor:\CAPETtmpnum}}%
     \label{capet:eno:\CAPETtmpnum}}%
  }
  % \CAPETANNEXE : rejoue tous les corriges captures, dans l'ordre, avec intertitres
  \newcommand{\CAPETANNEXE}{%
    \setcounter{CAPETannexe}{0}%
    \loop\ifnum\value{CAPETannexe}<\value{CAPETexo}%
      \stepcounter{CAPETannexe}%
      \edef\CAPETcurnum{\theCAPETannexe}%
      \ifcsname CAPETnotionstart\CAPETcurnum\endcsname
        \par\medskip
        {\normalsize\bfseries\color{colExo}\csuse{CAPETnotionstart\CAPETcurnum}}%
        \par\nobreak\smallskip
      \fi
      \begin{tcolorbox}[exoBox, breakable,
        title={Corrige de l'exercice \CAPETcurnum\
               \hfill\normalfont\footnotesize
               \hyperref[capet:eno:\CAPETcurnum]{enonce p.~\pageref{capet:eno:\CAPETcurnum}}}]%
      \label{capet:cor:\CAPETcurnum}%
      \csuse{CAPETcorrige\CAPETcurnum}%
      \end{tcolorbox}%
    \repeat
  }
\else
  % --- Modes historiques, comportement strictement inchange ---
  \newcommand{\CAPETSHOW}[2]{\ifdefined\CAPETCORRIGE #2\else #1\fi}
  \newcommand{\CAPETANNEXE}{}
\fi

% ── Difficulté 1 à 3 (étoiles pleines, amssymb) ───────────────────────
\newcommand{\difficulte}[1]{%
  \ifcase#1\relax
  \or $\bigstar$\;niveau 1/3%
  \or $\bigstar\bigstar$\;niveau 2/3%
  \or $\bigstar\bigstar\bigstar$\;niveau 3/3%
  \fi
}

% ── Style encadré extrait de document technique (banque d'exercices) ──
\tcbset{
  docTech/.style={
    breakable,
    enhanced,
    colback=colGrayLight,
    colframe=colSeparator,
    fonttitle=\bfseries\small\itshape,
    coltitle=black,
    colbacktitle=colGrayLight,
    left=4pt, right=4pt, top=3pt, bottom=3pt,
    before skip=4pt, after skip=4pt,
  }
}

% ── Style tcolorbox de section principale ────────────────────────────
% Refonte 2026 : bandeau titre coloré arrondi + fond très pâle, sans cadre.
\tcbset{
  ficheSection/.style={
    breakable, enhanced jigsaw,
    frame hidden, boxrule=0pt, arc=1.8mm,
    colback=#1!4!white,
    colbacktitle=#1,
    coltitle=white,
    fonttitle=\sffamily\bfseries,
    left=7pt, right=7pt, top=5pt, bottom=5pt,
    toptitle=3pt, bottomtitle=3pt,
    before skip=10pt, after skip=8pt,
  }
}

% ── Style encadré question type ──────────────────────────────────────
% Carte blanche, bandeau sombre, liseré très clair.
\tcbset{
  questiontype/.style={
    breakable, enhanced jigsaw,
    colback=white,
    frame hidden, boxrule=0pt, arc=1.8mm,
    colbacktitle=colElec!78!black,
    coltitle=white,
    fonttitle=\sffamily\bfseries\small,
    left=7pt, right=7pt, top=4pt, bottom=4pt,
    toptitle=2.5pt, bottomtitle=2.5pt,
    before skip=8pt, after skip=8pt,
    borderline={0.6pt}{0pt}{colElec!12!white},
  }
}

% ── Style encadré jury ───────────────────────────────────────────────
% Bandeau ambre, fond crème, sans cadre.
\tcbset{
  jury/.style={
    breakable, enhanced jigsaw,
    colback=colJuryBg,
    frame hidden, boxrule=0pt, arc=1.8mm,
    colbacktitle=colJuryFrame,
    coltitle=white,
    fonttitle=\sffamily\bfseries,
    left=7pt, right=7pt, top=4pt, bottom=4pt,
    toptitle=3pt, bottomtitle=3pt,
    before skip=10pt, after skip=8pt,
  }
}

% ── En-tête / pied de page communs (refonte 2026) ─────────────────────
% #1 code fiche, #2 titre court, #3 couleur domaine
\newcommand{\ficheHeaders}[3]{%
  \fancyhead[L]{\small\sffamily{\bfseries\textcolor{#3}{#1}}\hspace{0.7em}\textcolor{colSlate}{#2}}%
  \fancyhead[R]{\small\sffamily\textcolor{colSlate}{CAPET SII IE\, --- \,Session 2026}}%
  \fancyfoot[C]{\small\sffamily\textcolor{colSlate}{\thepage\,/\,\pageref{LastPage}}}%
  \fancyhead[C]{}\fancyfoot[L]{}\fancyfoot[R]{}%
}

% ── Bloc en-tête de fiche (BLOC 0, refonte 2026) ──────────────────────
% #1 code  #2 titre  #3 couleur domaine  #4 domaine  #5 ligne présence  #6 sous-thèmes
\newcommand{\ficheEntete}[6]{%
  \begin{tcolorbox}[enhanced, breakable,
    frame hidden, boxrule=0pt, arc=2.5mm,
    interior style={left color=#3, right color=#3!75!black},
    left=9pt, right=9pt, top=7pt, bottom=7pt,
    before skip=0pt, after skip=6pt]
    {\color{white}\sffamily\bfseries\LARGE #1\quad #2}\\[4pt]
    {\color{#3!22!white}\sffamily\small \MakeUppercase{#4}\quad\textbar\quad #5}\\[1pt]
    {\color{#3!35!white}\sffamily\footnotesize Sous-thèmes : #6}
  \end{tcolorbox}%
}

% ── Mise en page ─────────────────────────────────────────────────────
\pagestyle{fancy}
\fancyhf{}
\renewcommand{\headrulewidth}{0.5pt}
\renewcommand{\footrulewidth}{0pt}
\AtBeginDocument{%
  \renewcommand{\headrule}{{\color{colRule}\hrule width\headwidth height 0.5pt}}%
}

\setlist[itemize]{leftmargin=1.2em, itemsep=1pt, topsep=2pt, parsep=0pt}
\setlist[enumerate]{leftmargin=1.2em, itemsep=3pt, topsep=2pt, parsep=0pt}

\setlength{\parskip}{2pt}
\setlength{\columnsep}{10pt}

\def\CAPETBEGINFILE{\begin{document}}
\def\CAPETENDFILE{\end{document}}
\else
\def\CAPETBEGINFILE{}
\def\CAPETENDFILE{}
\fi

\ficheHeaders{T33}{Informatique embarquée}{colElec}

\CAPETBEGINFILE

% -----------------------------------------------------------------------
%  BLOC EN-TETE
% -----------------------------------------------------------------------

\ficheEntete{T33}{Informatique embarquée, programmation et réseaux industriels}{colElec}%
  {électronique numérique - Informatique et communication}%
  {\textbf{Présence :} \textbf{11/13} sujets (2013-2026) - thème \emph{très fréquent}}%
  {algorithmique et algorigramme - Python et C - opérations bit à bit - architecture d'un microcontrôleur - timers et interruptions - UART, RS-485, I2C, SPI - Modbus et Ethernet industriel - IHM - FPGA et VHDL}

% -----------------------------------------------------------------------
%  1. RAPPELS DE COURS
% -----------------------------------------------------------------------

\begin{tcolorbox}[
  ficheSection=colElec,
  title={1. Rappels de cours},
]

\textbf{1.1 Algorithmique et programmation}

\smallskip
Trois structures suffisent : la \textbf{séquence}, l'\textbf{alternative} (si\ldots alors\ldots sinon) et la \textbf{répétition} (pour, tant que). Un algorigramme les dessine : ovale début/fin, rectangle traitement, losange test, parallélogramme entrée/sortie.

\begin{center}
\renewcommand{\arraystretch}{1.3}
\begin{tabular}{@{}lll@{}}
\toprule
\textbf{Notion} & \textbf{Python} & \textbf{C} \\
\midrule
Boucle sur $n$ éléments & \texttt{for i in range(n):} (0 à $n-1$) & \texttt{for (i = 0; i < n; i++)} \\
Division entière / réelle & \texttt{7 // 2} $\rightarrow 3$ ; \texttt{7 / 2} $\rightarrow 3{,}5$ & \texttt{7 / 2} $\rightarrow 3$ entre entiers \\
Égalité / affectation & \texttt{==} / \texttt{=} & \texttt{==} / \texttt{=} \\
Types & dynamiques & déclarés (\texttt{int}, \texttt{float}, \texttt{uint8\_t}\ldots) \\
\bottomrule
\end{tabular}
\end{center}

\textbf{Opérations bit à bit} (registres, masques, trames) :
\begin{center}
\renewcommand{\arraystretch}{1.3}
\begin{tabular}{@{}lll@{}}
\toprule
\textbf{Action sur le bit de rang $k$} & \textbf{Écriture} & \textbf{Exemple ($k = 2$)} \\
\midrule
Tester & \texttt{(x >> k) \& 1}, ou \texttt{(x \& m) != 0} & masque \texttt{m = 0x04} \\
Mettre à 1 & \texttt{x |= (1 << k)} & \texttt{x |= 0x04} \\
Mettre à 0 & \texttt{x \&= \textasciitilde(1 << k)} & \texttt{x \&= 0xFB} \\
Inverser & \texttt{x \textasciicircum= (1 << k)} & \texttt{x \textasciicircum= 0x04} \\
\bottomrule
\end{tabular}
\end{center}
Sur $n$ bits : non signé de $0$ à $2^n - 1$ ; signé en complément à 2 de $-2^{n-1}$ à $2^{n-1}-1$ (sur 8 bits, \texttt{0xFF} vaut 255 ou $-1$).

\medskip\hrule\vspace{4pt}

\textbf{1.2 Architecture d'un microcontrôleur et composants programmables}

\smallskip
Un microcontrôleur réunit sur une puce un \textbf{processeur} (CPU), de la \textbf{mémoire} (Flash pour le programme, RAM pour les données, EEPROM pour les paramètres) et des \textbf{périphériques} reliés par des bus : ports d'entrées-sorties (GPIO), timers et MLI, CAN, liaisons série. Une mémoire de $n_{adr}$ bits d'adresse et de mots de $w$ bits contient $2^{n_{adr}} \times w$ bits.

\begin{itemize}
  \item \textbf{Scrutation ou interruption} : la scrutation teste périodiquement un état et peut manquer un événement bref ; une interruption suspend le programme pour exécuter une routine dès que l'événement survient.
  \item \textbf{Timer} : un compteur incrémenté à $f_{clk}/P$ ($P$ : prédiviseur) déborde après $ARR + 1$ coups, soit une période $T = P\,(ARR+1)/f_{clk}$ ; il sert de base de temps et génère la MLI.
  \item \textbf{FPGA} : circuit logique reconfigurable (tables de correspondance, bascules) décrit en \textbf{VHDL} ou Verilog. Une \texttt{entity} déclare les ports, une \texttt{architecture} décrit le comportement ; les instructions s'exécutent en parallèle, contrairement à un programme de microcontrôleur. Choix : FPGA pour le traitement rapide et parallèle, microcontrôleur pour la commande séquentielle à bas coût.
\end{itemize}

\medskip\hrule\vspace{4pt}

\textbf{1.3 Liaisons série et réseaux industriels}

\smallskip
En liaison série asynchrone (UART), chaque caractère est encadré par un bit de \textbf{start} (0) et un ou deux bits de \textbf{stop} (1), avec une parité éventuelle ; les données partent \textbf{bit de poids faible en premier}. En 8N1, un octet coûte 10 bits. Exemple : le caractère \og E \fg{} (\texttt{0x45}, 2023 E5).

\begin{center}
\begin{tikzpicture}[xscale=0.62, yscale=0.55, font=\scriptsize]
  % niveaux : repos 1, start 0, d0..d7 = 1 0 1 0 0 0 1 0, stop 1
  \draw[thick, colElec] (0,1) -- (1,1) -- (1,0) -- (2,0) -- (2,1) -- (3,1) -- (3,0) -- (4,0) -- (4,1) -- (5,1) -- (5,0) -- (8,0) -- (8,1) -- (9,1) -- (9,0) -- (10,0) -- (10,1) -- (12,1);
  \foreach \x in {1,...,11} \draw[gray!50, dashed] (\x,-0.2) -- (\x,1.3);
  \node at (1.5,1.6) {start};
  \foreach \i/\b in {0/1,1/0,2/1,3/0,4/0,5/0,6/1,7/0} {
    \node at (2.5+\i,1.6) {d\i};
    \node[colElec!70!black] at (2.5+\i,-0.6) {\b};
  }
  \node at (10.5,1.6) {stop};
  \node[anchor=east] at (-0.1,1) {repos = 1};
  \node[anchor=west] at (12.2,0.5) {«~E~» = \texttt{0x45}, 8N1};
\end{tikzpicture}
\end{center}

\begin{center}
\renewcommand{\arraystretch}{1.3}
\small
\begin{tabular}{@{}p{2.5cm}p{3.6cm}p{3.0cm}p{5.6cm}@{}}
\toprule
\textbf{Liaison / réseau} & \textbf{Support} & \textbf{Topologie} & \textbf{Repères} \\
\midrule
RS-232 & tensions $\pm$ 3 à 15\,V & point à point & 15\,m, débits faibles \\
RS-485 & paire différentielle & multipoint & 32 nœuds, 1\,200\,m \\
I2C & SDA, SCL (collecteur ouvert) & bus, adresses 7 bits & résistances de tirage, 100/400\,kbit/s \\
SPI & MOSI, MISO, SCK, CS & maître unique & full-duplex, un CS par esclave \\
Bus CAN & paire différentielle & multimaître & arbitrage, voir T25 \\
Modbus RTU & sur RS-485 & maître / esclaves & adresses 1 à 247, CRC16, silence 3{,}5 caractères \\
Modbus TCP, Profinet & Ethernet & étoile (commutateurs) & port 502 pour Modbus TCP \\
\bottomrule
\end{tabular}
\end{center}

\medskip\hrule\vspace{4pt}

\textbf{1.4 IHM et supervision}

\smallskip
L'\textbf{interface homme-machine} affiche l'état du système, permet de commander et de régler, et gère les \textbf{alarmes} (horodatées, acquittées). Elle dialogue avec l'automate (API) par un réseau (Modbus TCP, Profinet) en lisant et écrivant des variables. Une supervision (SCADA) centralise plusieurs installations. Règles d'ergonomie : couleurs normalisées doublées d'un texte ou d'un symbole, informations critiques visibles sans navigation, confirmation des commandes dangereuses.

\end{tcolorbox}

% -----------------------------------------------------------------------
%  2. FORMULAIRE
% -----------------------------------------------------------------------

\begin{tcolorbox}[
  ficheSection=colElec,
  title={2. Formulaire synthétique},
]
\small
\begin{multicols}{2}
\renewcommand{\arraystretch}{1.45}
\begin{tabular}{@{}p{4.4cm}ll@{}}
\toprule
\textbf{Formule} & \textbf{Usage} & \textbf{U.} \\
\midrule
\multicolumn{3}{@{}l}{\footnotesize\textit{Codage}} \\
$0 \le x \le 2^n - 1$ & Non signé & - \\
$-2^{n-1} \le x \le 2^{n-1} - 1$ & Complément à 2 & - \\
$C = 2^{n_{adr}} \times w$ & Mémoire & bit \\
\midrule
\multicolumn{3}{@{}l}{\footnotesize\textit{Temps}} \\
$t = N_{cycles}/f_{clk}$ & Exécution & s \\
$T = P\,(ARR+1)/f_{clk}$ & Timer & s \\
$E \approx \sum_i \phi_i\,\Delta t$ & Rectangles & J \\
\bottomrule
\end{tabular}

\columnbreak

\renewcommand{\arraystretch}{1.45}
\begin{tabular}{@{}p{4.4cm}ll@{}}
\toprule
\textbf{Formule} & \textbf{Usage} & \textbf{U.} \\
\midrule
\multicolumn{3}{@{}l}{\footnotesize\textit{Liaisons série}} \\
$t_{car} = (1 + n_d + n_p + n_s)/R$ & Caractère & s \\
$D_u = R\,n_d/(1 + n_d + n_p + n_s)$ & Débit utile & bit/s \\
$t_{trame} = N\,t_{car}$ & Trame & s \\
$t_{sil} \ge 3{,}5\,t_{car}$ & Modbus RTU & s \\
$N_{adr} = 2^7 = 128$ & I2C & - \\
\bottomrule
\end{tabular}
\end{multicols}
\normalsize
\end{tcolorbox}

% -----------------------------------------------------------------------
%  3. QUESTIONS TYPES
% -----------------------------------------------------------------------

\begin{tcolorbox}[
  ficheSection=colElec,
  title={3. Questions types - extraites des sujets réels},
]

\begin{tcolorbox}[questiontype,
  title={QT1 - Lire un bit d'alarme avec un masque \hfill\normalfont\footnotesize E1 2026 (D10)}
]
\textbf{Énoncé.}
Pour programmer le calculateur du bus, on veut accéder directement à l'état de l'alerte \og température pile \fg{} à partir d'une variable \texttt{Datafield} de type entier long. Proposer une fonction en pseudo-langage qui renvoie 1 en cas d'alerte et 0 sinon.

\tcbline

\textbf{Corrigé.} \textit{(corrigé officiel : rapport 2026 du troisième concours, question D11)}

\textbf{Étape 1 - Masque.} Le bit d'alerte est le bit de rang 2 de \texttt{Datafield} ; le masque ne garde que lui : \texttt{Masque = 0004} (hexadécimal).

\textbf{Étape 2 - Fonction.}
\begin{center}
\begin{tabular}{@{}l@{}}
\texttt{Si (Datafield ET Masque) = 0 alors renvoyer 0} \\
\texttt{Sinon renvoyer 1} \\[3pt]
Python : \texttt{return 1 if datafield \& 0x0004 else 0} \\
C : \texttt{return (datafield \& 0x0004) != 0;}
\end{tabular}
\end{center}

\textbf{Conclusion.} Le ET bit à bit isole un seul bit sans dépendre des autres alarmes de l'octet ; c'est la méthode attendue, plus sûre qu'une comparaison de l'octet entier à une valeur.
\end{tcolorbox}

\vspace{4pt}

\begin{tcolorbox}[questiontype,
  title={QT2 - Durée d'une trame GPS en liaison série \hfill\normalfont\footnotesize E2 disciplinaire IE 2023 (E4-E8)}
]
\textbf{Énoncé.}
Un module GPS transmet des caractères ASCII sur une liaison série à \SI{115200}{bauds}, avec 1 bit de start, 8 bits de données et 1 bit de stop. Une trame compte 71 caractères. Le robot roule au plus à \SI{6}{km/h} et sa position doit être connue à \SI{1}{cm} près. (a) Combien de bits par caractère ? (b) Quelle est la valeur transmise pour \og E \fg{} ? (c) Durée d'une trame ? (d) Distance parcourue pendant ce temps ; conclure.

\tcbline

\textbf{Corrigé.} \textit{(d'après le corrigé officiel, rapport 2023)}

\textbf{Étape 1 - Bits par caractère.} $1 + 8 + 1 = \mathbf{10}$ bits. \og E \fg{} $= \mathtt{0x45} = 0100\,0101_2$.

\textbf{Étape 2 - Durée de la trame.}
\[
  t_{trame} = \frac{71 \times 10}{115\,200} = \mathbf{6{,}16\ ms}
\]

\textbf{Étape 3 - Distance parcourue.}
\[
  d = v\,t = \frac{6}{3{,}6} \times 6{,}16\cdot10^{-3} = \mathbf{1{,}03\ cm}
\]

\textbf{Point de vigilance.} Le rapport écrit \og 1{,}03 m \fg{} : l'unité est fausse, le calcul donne 1{,}03 cm. La conclusion officielle reste juste.

\textbf{Conclusion.} 1{,}03\,cm dépasse la précision exigée de 1\,cm, de peu : la liaison est trop lente. Il faut un débit plus élevé ou une trame plus courte.
\end{tcolorbox}

\vspace{4pt}

\begin{tcolorbox}[questiontype,
  title={QT3 - Intégrer une irradiance et compléter l'algorithme \hfill\normalfont\footnotesize E2 disciplinaire IE 2025 (Q4-Q5)}
]
\textbf{Énoncé.}
Un relevé d'irradiance $\phi$ (W/m$^2$) comporte 3 points par minute. (a) Préciser le pas et calculer, par la méthode des rectangles, l'énergie surfacique reçue en 3 minutes sachant que la somme des 9 valeurs vaut 1\,327\,W/m$^2$. (b) Écrire l'algorithme qui calcule l'énergie annuelle en kWh/m$^2$ à partir d'un tableau \texttt{Data} (temps en s, irradiance) relevé toutes les heures.

\tcbline

\textbf{Corrigé.} \textit{(d'après le corrigé officiel, rapport 2025)}

\textbf{Étape 1 - Pas et énergie.} Trois points par minute : $\text{Pas} = \SI{20}{s}$.
\[
  E_s = \sum_i \phi_i\,\text{Pas} = 1\,327 \times 20 = 26\,540\ \mathrm{J/m^2} = \mathbf{7{,}37\ Wh/m^2}
\]

\textbf{Étape 2 - Algorithme} (version Python) :
\begin{center}
\begin{tabular}{@{}l@{}}
\texttt{irrad = 0} \\
\texttt{for i in range(len(Data) - 1):} \\
\texttt{\ \ \ \ dt = Data[i+1][0] - Data[i][0] \ \ \# en s} \\
\texttt{\ \ \ \ irrad = irrad + Data[i][1] * dt \ \# J/m2} \\
\texttt{print(irrad / 3.6e6, "kWh/m2")}
\end{tabular}
\end{center}

\textbf{Point de vigilance.} Le corrigé officiel boucle jusqu'à la dernière ligne alors qu'il lit l'indice suivant, et divise par 1\,000 : c'est juste si l'écart de temps est en heures ; avec un temps en secondes, il faut diviser par $3{,}6\cdot10^6$.

\textbf{Conclusion.} La méthode des rectangles suffit pour une irradiance relevée à pas fixe ; le résultat annuel (1\,447\,kWh/m$^2$ dans le sujet) sert ensuite à dimensionner la surface de panneaux.
\end{tcolorbox}

\end{tcolorbox}

% -----------------------------------------------------------------------
%  4. ATTENTES DU JURY
% -----------------------------------------------------------------------

\begin{tcolorbox}[
  jury,
  title={4. Attentes du jury - Erreurs récurrentes et points de vigilance},
]

\begin{enumerate}[label=\textbf{\arabic*.}, leftmargin=*, itemsep=5pt]
  \item \textbf{Compter 8 bits par caractère.} En liaison asynchrone, ajouter start, stop et parité : 8N1 fait 10 bits (2023 E4, 2018).
  \item \textbf{Lire une trame série bit de poids fort en premier.} L'UART émet le bit de poids faible juste après le start.
  \item \textbf{Confondre bauds et débit utile.} À 115\,200 bauds en 8N1, on transmet 11\,520 octets de données par seconde.
  \item \textbf{Mal borner une boucle.} Indices de 0 à $n-1$ ; lire \texttt{t[i+1]} impose de s'arrêter à $n-2$.
  \item \textbf{Oublier les parenthèses d'un masque.} En C, \texttt{==} est prioritaire sur \texttt{\&} : écrire \texttt{(x \& m) == 0}.
  \item \textbf{Ignorer le signe d'une variable.} Sur 8 bits signés, \texttt{0xFF} vaut $-1$.
  \item \textbf{Confondre les liaisons.} RS-232 point à point ; RS-485 multipoint différentielle ; I2C à résistances de tirage ; SPI avec un fil de sélection par esclave.
  \item \textbf{Manquer un événement bref par scrutation.} Utiliser une interruption.
\end{enumerate}

\medskip
% Rapport_Jury_2018.md l.158 (avant-propos, commun à toutes les options) - verbatim vérifié
\begin{tcolorbox}[
  colback=colJuryFrame!12!colJury,
  colframe=colJuryFrame,
  left=4pt, right=4pt, top=2pt, bottom=2pt,
  fontupper=\small\itshape
]
«~Depuis la session 2017, que ce soit à l’agrégation ou au CAPET, à l’externe comme à l’interne, il existe une option ingénierie informatique et surtout un programme d’informatique pour toutes les options de ces concours. Les compétences liées à ce programme sont par conséquent également évaluées. Le jury conseille aux futurs candidats de ne surtout pas faire l’impasse sur l’informatique [\ldots]~»\\
\upshape\footnotesize Rapport du jury 2018, avant-propos (commun à toutes les options).
\end{tcolorbox}

\end{tcolorbox}

% -----------------------------------------------------------------------
%  5. PRÉSENCE
% -----------------------------------------------------------------------

\begin{tcolorbox}[
  ficheSection=colPresence,
  title={Présence dans les sujets analysés (2013-2026)},
]
\begin{center}
\footnotesize\renewcommand{\arraystretch}{1.15}
\begin{tabular}{lccccccccccccc}
\toprule
\textbf{Année} & 13 & 14 & 15 & 16 & 17 & 18 & 19 & 20 & 21 & 22 & 23 & 24 & 25 \\
\midrule
Épreuve commune & $\checkmark$ & $\checkmark$ & - & - & $\checkmark$ & - & $\checkmark$ & $\checkmark$ & $\checkmark$ & $\checkmark$ & $\checkmark$ & $\checkmark$ & $\checkmark$ \\
E2 discipl. IE  & - & - & - & - & $\checkmark$ & $\checkmark$ & $\checkmark$ & $\checkmark$ & $\checkmark$ & $\checkmark$ & $\checkmark$ & $\checkmark$ & $\checkmark$ \\
\bottomrule
\end{tabular}
\end{center}
\smallskip
\footnotesize\textbf{Sous-thèmes rencontrés :} algorithme ou algorigramme à compléter (2014, 2017, 2019--2025) - Python (2019, 2021, 2022, 2023 ; E2 2026) - langage C (2020) - liaison série et ASCII (2018, 2023, 2025) - Modbus et Ethernet (2013, 2014, 2023, 2025) - IHM (2018) - I2C (2022) - pseudo-langage avec masque (E1 2026).
\end{tcolorbox}

\CAPETENDFILE
